%%% This file is automatically generated by LeanArchitect. %%%

%%% Macro definitions for \inputleannode, \inputleanmodule %%%

\makeatletter

% \newleannode{name}{latex} defines a new Lean node
\providecommand{\newleannode}[2]{%
  \expandafter\gdef\csname leannode@#1\endcsname{#2}}
% \inputleannode{name} inputs a Lean node
\providecommand{\inputleannode}[1]{%
  \csname leannode@#1\endcsname}

% \newleanmodule{module}{latex} defines a new Lean module
\providecommand{\newleanmodule}[2]{%
  \expandafter\gdef\csname leanmodule@#1\endcsname{#2}}
% \inputleanmodule{module} inputs a Lean module
\providecommand{\inputleanmodule}[1]{%
  \csname leanmodule@#1\endcsname}

\makeatother

%%% Start of main content %%%

\newleannode{everyDef}{\begin{definition}[Dirichlet nonvanishing premise]
\label{everyDef}
\lean{UpperTheory.EveryCharacter}
% at /home/ubuntu/riemann-original-apolge/leanarchitect-blueprint/RiemannArchitect.lean:8.0-11.130
\leanok
Every Dirichlet L-function is nonzero to the right of one half, excluding the principal-character point at one.
\end{definition}
}

\newleannode{half}{\begin{definition}[Right half-plane predicate]
\label{half}
\lean{Claim.HalfPlane}
% at /home/ubuntu/riemann-original-apolge/leanarchitect-blueprint/RiemannArchitect.lean:14.0-17.130
\leanok
Nonvanishing in the open right half-plane for every positive level, subject to the exception at one.
\end{definition}
}

\newleannode{stripDef}{\begin{definition}[Strip predicate]
\label{stripDef}
\lean{Claim.Strip}
% at /home/ubuntu/riemann-original-apolge/leanarchitect-blueprint/RiemannArchitect.lean:20.0-23.134
\leanok
A Dirichlet L-function zero in the open critical strip has real part one half.
\end{definition}
}

\newleannode{zetaDef}{\begin{definition}[Zeta-zero predicate]
\label{zetaDef}
\lean{Claim.Zeta}
% at /home/ubuntu/riemann-original-apolge/leanarchitect-blueprint/RiemannArchitect.lean:26.0-29.83
\leanok
Every nontrivial zeta zero distinct from one lies on the critical line.
\end{definition}
}

\newleannode{factor}{\begin{theorem}[Euler factor is nonzero]
\label{factor}
\lean{Claim.factor_ne_zero}
% at /home/ubuntu/riemann-original-apolge/leanarchitect-blueprint/RiemannArchitect.lean:33.0-49.22
\leanok
At positive real part the prime-power term has norm less than one, so a deleted Euler factor cannot vanish.
\end{theorem}
\begin{proof}
\leanok
\end{proof}
}

\newleannode{primZero}{\begin{theorem}[Pass a zero to the primitive character]
\label{primZero}
\lean{Claim.prim_zero}
% at /home/ubuntu/riemann-original-apolge/leanarchitect-blueprint/RiemannArchitect.lean:54.0-66.80
\leanok
Change level to the conductor and rule out zeros of the deleted Euler factors.
\end{theorem}
\begin{proof}
\uses{factor}
\leanok
\end{proof}
}

\newleannode{gamma}{\begin{theorem}[Gamma factor is nonzero]
\label{gamma}
\lean{Claim.gamma_ne_zero}
% at /home/ubuntu/riemann-original-apolge/leanarchitect-blueprint/RiemannArchitect.lean:71.0-78.93
\leanok
Both possible archimedean factors are nonzero at positive real part.
\end{theorem}
\begin{proof}
\leanok
\end{proof}
}

\newleannode{reflect}{\begin{theorem}[Reflect a primitive zero]
\label{reflect}
\lean{Claim.prim_reflect}
% at /home/ubuntu/riemann-original-apolge/leanarchitect-blueprint/RiemannArchitect.lean:83.0-100.13
\leanok
The completed functional equation transfers a primitive zero to the inverse character at one minus the original argument.
\end{theorem}
\begin{proof}
\uses{gamma}
\leanok
\end{proof}
}

\newleannode{strip}{\begin{theorem}[All Dirichlet zeros in the strip]
\uses{half,stripDef}
\label{strip}
\lean{Claim.strip}
% at /home/ubuntu/riemann-original-apolge/leanarchitect-blueprint/RiemannArchitect.lean:105.0-121.45
\leanok
Use half-plane nonvanishing on the right. On the left, pass to the primitive character and reflect its zero to the right.
\end{theorem}
\begin{proof}
\uses{primZero,reflect}
\leanok
\end{proof}
}

\newleannode{zeta}{\begin{theorem}[Zeta consequence of nonvanishing]
\uses{half,zetaDef}
\label{zeta}
\lean{Claim.zeta}
% at /home/ubuntu/riemann-original-apolge/leanarchitect-blueprint/RiemannArchitect.lean:126.0-153.97
\leanok
Use the level-one character in the strip and the zeta functional equation and nonvanishing outside it.
\end{theorem}
\begin{proof}
\uses{strip}
\leanok
\end{proof}
}

\newleannode{right}{\begin{theorem}[Reflect a hypothetical zeta zero]
\label{right}
\lean{Formalism.VW.rh_of_right}
% at /home/ubuntu/riemann-original-apolge/leanarchitect-blueprint/RiemannArchitect.lean:157.0-216.34
\leanok
Right-half-plane nonvanishing implies RH by the zeta functional equation, with the exceptional gamma and trigonometric factors handled explicitly.
\end{theorem}
\begin{proof}
\leanok
\end{proof}
}

\newleannode{vGamma}{\begin{theorem}[Archimedean factor for the primitive proof]
\label{vGamma}
\lean{Formalism.VW.gammaFactor_ne_zero}
% at /home/ubuntu/riemann-original-apolge/leanarchitect-blueprint/RiemannArchitect.lean:219.0-224.93
\leanok
Compute the gamma factor according to character parity on the positive half-plane.
\end{theorem}
\begin{proof}
\leanok
\end{proof}
}

\newleannode{crit}{\begin{theorem}[Primitive zeros on the critical line]
\uses{everyDef}
\label{crit}
\lean{Formalism.VW.critLine}
% at /home/ubuntu/riemann-original-apolge/leanarchitect-blueprint/RiemannArchitect.lean:227.0-245.37
\leanok
Half-plane nonvanishing and the primitive functional equation force a zero in the strip onto the critical line.
\end{theorem}
\begin{proof}
\uses{vGamma}
\leanok
\end{proof}
}

\newleannode{rhEvery}{\begin{theorem}[Specialize to the level-one character]
\uses{everyDef}
\label{rhEvery}
\lean{Formalism.VW.rh_of_every}
% at /home/ubuntu/riemann-original-apolge/leanarchitect-blueprint/RiemannArchitect.lean:248.0-253.37
\leanok
Apply the zeta reduction to the trivial character modulo one, whose L-function is zeta.
\end{theorem}
\begin{proof}
\uses{right}
\leanok
\end{proof}
}

\newleanmodule{RiemannArchitect}{
\inputleannode{everyDef}

\inputleannode{half}

\inputleannode{stripDef}

\inputleannode{zetaDef}

\inputleannode{factor}

\inputleannode{primZero}

\inputleannode{gamma}

\inputleannode{reflect}

\inputleannode{strip}

\inputleannode{zeta}

\inputleannode{right}

\inputleannode{vGamma}

\inputleannode{crit}

\inputleannode{rhEvery}

\begin{theorem}[Primitive zero-free criterion]
\label{topG}
Combine the top-level condition with the required host data for each primitive character.
\end{theorem}
\begin{proof}
Source-traced declaration \texttt{UpperTheory.Top.zeroFreeAll\_of\_topG}.
\end{proof}


\begin{theorem}[Primitive character nonvanishing]
\label{greenZero}
The Green and window hypotheses and a definite-side witness provide the inputs of the primitive criterion.
\end{theorem}
\begin{proof}
\uses{topG}
Source-traced declaration \texttt{UpperTheory.Top.zero\_free\_all\_of\_green\_v\_win}.
\end{proof}


\begin{theorem}[Extend to imprimitive characters]
\label{imprimitive}
Use the character of minimal conductor and show that the remaining Euler factors cannot vanish.
\end{theorem}
\begin{proof}
Source-traced declaration \texttt{UpperTheory.Top.imprimitive}.
\end{proof}


\begin{theorem}[From windows to all characters]
\label{greenEvery}
First derive the primitive zero-free result, then extend to arbitrary positive levels.
\end{theorem}
\begin{proof}
\uses{greenZero,imprimitive}
Source-traced declaration \texttt{UpperTheory.Top.every\_character\_of\_green\_v\_win}.
\end{proof}


\begin{theorem}[Construct a definite-side witness]
\label{definite}
Theta integrals, local matching, a covolume bound, and local parameters supply a nonempty family of definite-side hypotheses.
\end{theorem}
\begin{proof}
Source-traced declaration \texttt{UpperTheory.Lit.nonempty\_hyps\_definite\_side\_of\_closed\_bound}.
\end{proof}


\begin{theorem}[Theta integrals and local matching]
\label{match}
Combine the definite-side witness with the Green and window input to obtain the all-character conclusion.
\end{theorem}
\begin{proof}
\uses{definite,greenEvery}
Source-traced declaration \texttt{UpperTheory.Lit.every\_character\_of\_theta\_integrals\_local\_matching}.
\end{proof}


\begin{theorem}[Assemble the literature inputs]
\label{literature}
Package published results, local matching, theta integrals and a covolume estimate for the theta-to-nonvanishing theorem.
\end{theorem}
\begin{proof}
\uses{match}
Source-traced declaration \texttt{UpperTheory.LiteratureTheorems.every\_character\_of\_literature\_of\_vWin}.
\end{proof}


\begin{theorem}[Galois row for the local parameter]
\label{row}
The source obtains its local–global row by applying a Galois-representations result to a theta lift.
\end{theorem}
\begin{proof}
Source-traced declaration \texttt{UpperTheory.Gal.Int.galoisLocalGlobalStrong\_recParameter}.
\end{proof}


\begin{theorem}[Local parameter construction]
\label{local}
The project constructs the local parameter with the required compatibilities.
\end{theorem}
\begin{proof}
Source-traced declaration \texttt{UpperTheory.Zero.LLC.recTExists}.
\end{proof}


\begin{theorem}[Window condition from a row]
\label{window}
Use open-image and base-change inputs together with the given row to obtain the automorphic window.
\end{theorem}
\begin{proof}
Source-traced declaration \texttt{UpperTheory.One.vwin\_of\_row}.
\end{proof}


\begin{theorem}[Nonvanishing given the Galois row]
\label{ofRow}
Assemble the literature package using the supplied row and local parameter existence. Apply the window result to that same row.
\end{theorem}
\begin{proof}
\uses{literature,local,window}
Source-traced declaration \texttt{UpperTheory.FinalTheorems.everyCharacter\_of\_row}.
\end{proof}


\begin{theorem}[Nonvanishing for all characters]
\label{every}
Discharge the row assumption with the local–global result. This is the unconditional input consumed by the public analytic deductions.
\end{theorem}
\begin{proof}
\uses{ofRow,row}
Source-traced declaration \texttt{UpperTheory.FinalTheorems.everyCharacter}.
\end{proof}


\begin{theorem}[Public nonvanishing bridge]
\label{allChi}
Expose the unconditional all-character result to the final deduction module.
\end{theorem}
\begin{proof}
\uses{every}
Source-traced declaration \texttt{allChi}.
\end{proof}


\begin{theorem}[Riemann hypothesis from all characters]
\label{rhAll}
Apply the checked level-one reduction to the nonvanishing bridge.
\end{theorem}
\begin{proof}
\uses{rhEvery,allChi}
Source-traced declaration \texttt{rhOfAll}.
\end{proof}


\begin{theorem}[Primitive critical-line bridge]
\label{critAll}
Apply the primitive critical-line theorem with the unconditional nonvanishing input.
\end{theorem}
\begin{proof}
\uses{crit,allChi}
Source-traced declaration \texttt{critAll}.
\end{proof}


\begin{theorem}[Package the three analytic claims]
\label{claimAll}
Assemble the half-plane result, the general Dirichlet strip result and the zeta statement in a conjunction.
\end{theorem}
\begin{proof}
\uses{allChi,strip,zeta}
Source-traced declaration \texttt{claimAll}.
\end{proof}


\begin{theorem}[Riemann hypothesis]
\label{mainRH}
The first public theorem takes its proof from the zeta bridge.
\end{theorem}
\begin{proof}
\uses{rhAll}
Source-traced declaration \texttt{riemannHypothesis}.
\end{proof}


\begin{theorem}[Nontrivial zeta zeros on the line]
\label{mainZeta}
The second public statement uses the same bridge under a definitionally equal type.
\end{theorem}
\begin{proof}
\uses{rhAll}
Source-traced declaration \texttt{riemannZeta\_eq\_zero\_re}.
\end{proof}


\begin{theorem}[Dirichlet half-plane nonvanishing]
\label{mainNV}
The public theorem is the all-character nonvanishing bridge.
\end{theorem}
\begin{proof}
\uses{allChi}
Source-traced declaration \texttt{LFunction\_ne\_zero\_of\_half\_lt\_re}.
\end{proof}


\begin{theorem}[All Dirichlet zeros on the line]
\label{mainStrip}
Project the strip proof from the three-part conjunction. Only the strip component is selected at this use.
\end{theorem}
\begin{proof}
\uses{claimAll}
Source-traced declaration \texttt{LFunction\_eq\_zero\_re}.
\end{proof}


\begin{theorem}[Primitive Dirichlet zeros on the line]
\label{mainPrim}
The original public theorem uses the dedicated primitive bridge.
\end{theorem}
\begin{proof}
\uses{critAll}
Source-traced declaration \texttt{LFunction\_eq\_zero\_re\_of\_isPrimitive}.
\end{proof}

}